% Five Lagrange points in the small-secondary-mass limit.
% Reference: handdraw/lagrange_point.png. Compile: pdflatex lagrange_point_tikz.tex
\documentclass[tikz,border=6pt]{standalone}
\usepackage{amsmath}
\usepackage{fix-cm}
\usetikzlibrary{arrows.meta}
\begin{document}
\begin{tikzpicture}[>=Stealth,thick,every node/.style={font=\fontsize{18}{22}\selectfont}]
  \draw[gray] (-3.3,0) -- (5,0);
  \draw[blue] (0,0) circle (.77);
  \draw[blue] (3.4,0) circle (.38);
  \fill[red] (0,0) circle (2.2pt);
  \fill[red] (3.4,0) circle (2.2pt);
  % L4,L5 form equilateral triangles with the two body centres.
  \foreach \x/\y/\name in {2.5/0/1,4.3/0/2,-2.8/0/3,1.7/2.9445/4,1.7/-2.9445/5} {
    \fill[gray] (\x,\y) ++(90:.11) -- ++(-60:.21) -- ++(180:.21) -- cycle;
  }
  \node[gray,below] at (2.5,-.2) {$L_1$};
  \node[gray,below] at (4.3,-.2) {$L_2$};
  \node[gray,below] at (-2.8,-.2) {$L_3$};
  \node[gray,right] at (1.9,2.9445) {$L_4$};
  \node[gray,right] at (1.9,-2.9445) {$L_5$};
  \node[red,anchor=west] at (-3.25,-2.2) {$m_1\gg m_2$};
\end{tikzpicture}
\end{document}
